\input{preamble}

\begin{document}

\title{Symplectic geometry}
\maketitle

\phantomsection
\label{section-phantom}

\tableofcontents

\noindent
See
\href{http://github.com/danimalabares/sg}%
{github.com/danimalabares/sg}
for the complete notes of Henrique Bursztyn's
course.

\section{Symplectic vector spaces}
\label{section-symplectic-vector-spaces}

\begin{proposition}
\label{proposition-symplectic-basis}
Every symplectic vector space $(V,\omega)$
admits a {\it symplectic basis},
i.e. a basis $e_1,…,e_n,f_1,…,f_n$
such that
$$
\omega(e_i,e_j)=0,\qquad
\omega(e_i,f_j)=1,\qquad \omega(f_i,f_j)=0.
$$
\end{proposition}

\begin{proof}
Like Gram-Schmidt, start with any vector
$e_1$,
then there exists another vector $f_1$
such that $\omega(e_1,f_1)=1$.
The space generated by $e_1,f_1$
is symplectic, so we get a normal
decomposition…
\end{proof}

\begin{exercise}
\label{exercise-omega-power-n}
Let $V$ be a vector space and $\omega$ a
2-form on $V$.
Show that $\omega$ is nondegenerate if and
only if
$\omega^n$ is not zero.
\end{exercise}

\begin{proof}
\begin{reference}
\href{https://math.stackexchange.com/
questions/417365/
how-to-show-omegan-is-a-volume-form-in-a-
symplectic-manifold-m-omega}
{Stack Exchange}
\end{reference}
If $\omega$ is nondegenerate,
by Proposition
\ref{proposition-symplectic-basis}
there is a symplectic basis
$e_1,…,e_n,f_1,…,f_n$
such that $\omega=\sum e^i \wedge f^i$.
Then
$$
\omega^n=n! e^1\wedge…\wedge e^n \wedge
f^1\wedge… \wedge f^n.
$$
Hence evaluating in
$e_1,…,e_n,f_1,…,f_n$
we see $\omega^n\neq 0$.

For the converse we suppose that $\omega$
is degenerate and prove that $\omega^n=0$.
It follows from the following property of
interior product:
$$
i_v\omega^n=n i_v\omega \wedge
\omega^{n-1}\qquad \forall  v \in V.
$$
Supposing that $\omega$ is degenerate we
have $v \in V$ such that $i_v\omega=0$,
which gives $i_v \omega^n=0$.
Completing $v$ to a basis $v_2,…,v_n$
we see that
$$
\omega^n(v,v_2,…,v_n)
=i_v\omega^n(v_2,…,v_n)=0.
$$
\end{proof}

\section{Hamiltonian vector fields, Poisson
brackets}
\label{section-hamil-vecto-field-poiss}

\noindent
We know that a symplectic form on a manifold
$M$ gives us a map
\begin{align*}
\omega^\sharp: TM &\longrightarrow T^*M \\
X &\longmapsto i_X\omega
\end{align*}

\noindent
which says that for $f \in C^\infty(M)$,
$$
i_{X_f}\omega=df.
$$
\begin{definition}
\label{definition-poisson-bracket}
\begin{align*}
\{\cdot,\cdot\}: C^\infty(M) \times
C^\infty(M) &\longrightarrow C^\infty(M) \\
\{f,g\}&=\omega(X_f,X_g)=df(X_g)=\mathcal{L}
_{X_f}g=-\mathcal{L}_{X_g}f
\end{align*}
is the {\it Poisson bracket}.
\end{definition}

\section{Symplectic structure of the
cotangent bundle}
\label{section-sympl-struc-cotan-bundl}

\noindent
Upshot. A tangent vector to $T^*M$
records two infinitesimal motions: the base
point moves in $M$, and the covector moves in
the cotangent fiber. The tautological 1-form
ignores this second motion. It takes the
covector which is the point of $T^*M$ and
evaluates it on the motion of the base point
using the differential of the projection
map of the bundle.

\medskip\noindent
Let $\pi:T^*M \to M$ be the cotangent bundle.
The {\it tautological 1-form} $\lambda \in
\Lambda^1(T^*M)$
is defined by
$$
\lambda_\xi(v)=\xi(d\pi_\xi(v)),\qquad \xi
\in T^*M,\; v \in T_\xi T^*M.
$$
The canonical symplectic form on $T^*M$ is
$$
\omega_{\text{can}}=-d\lambda.
$$
If $(q^i)$ are coordinates on $M$
and $(q^i,p_i)$ are the induced coordinates
on $T^*M$, then
$$
\lambda=\sum_i p_i dq^i,\qquad
\omega_{\text{can}}=\sum_i dq^i \wedge dp_i.
$$

\begin{exercise}
\label{exercise-cotan-bundl-canon-form}
Let $\alpha \in \Lambda^1(M)$ and view it as
a section
$\alpha:M \to T^*M$.
\begin{enumerate}
\item Show that $\alpha^*\lambda=\alpha$.
\item Conclude that
$\alpha^*\omega_{\text{can}}=-d\alpha$.
\item Show that the image $\alpha(M) \subset
T^*M$ is Lagrangian
if and only if $\alpha$ is closed.
\item In particular, show that the zero
section is Lagrangian.
\end{enumerate}
\end{exercise}

\begin{proof}
\begin{enumerate}
\item Pick $v \in TM$.
$$
\alpha^*\lambda(v)=\lambda(\alpha_*(v))
=\xi(\pi_*(\alpha_*(v))
$$
where $\alpha_*(v) \in T_\xi(M)$.
Now since $\alpha$ is a section,
$\pi \circ \alpha=\text{id}$,
and thinking about it, $\xi=\alpha$!
\end{enumerate}
\end{proof}

\begin{definition}
\label{definition-cotangent-lift}
Let $\varphi:M \to N$ be a diffeomorphism.
The {\it cotangent lift} of $\varphi$ is
the map
$$
\widehat{\varphi}:T^*M \longrightarrow T^*N
$$
defined as follows.
If $\xi \in T_p^*M$, then
$\widehat{\varphi}(\xi)$ is the covector at
$\varphi(p)$ given by
$$
\widehat{\varphi}(\xi)(w)
=\xi(d\varphi_p^{-1}(w)),
\qquad w \in T_{\varphi(p)}N.
$$
\end{definition}

\noindent
In other words, the cotangent lift moves
the base point by $\varphi$ and transports
the covector by the inverse differential.

\begin{proposition}
\label{proposition-cotangent-lift-symplectic}
The cotangent lift preserves the tautological
1-form:
$$
\widehat{\varphi}^*\lambda_N=\lambda_M.
$$
In particular,
$$
\widehat{\varphi}^*\omega_{\text{can},N}
=\omega_{\text{can},M},
$$
so cotangent lifts are symplectomorphisms.
\end{proposition}

\begin{proof}
Let $\xi \in T_p^*M$ and
$v \in T_\xi T^*M$.
The cotangent lift covers $\varphi$, so the
diagram
$$
\xymatrix{
T^*M \ar[r]^{\widehat{\varphi}}
\ar[d]_{\pi_M}
& T^*N \ar[d]^{\pi_N} \\
M \ar[r]_{\varphi} & N
}
$$
commutes.
By definition of the tautological 1-form,
$$
(\widehat{\varphi}^*\lambda_N)_\xi(v)
=\widehat{\varphi}(\xi)
\left(
d\pi_N(d\widehat{\varphi}_\xi(v))
\right).
$$
By commutativity of the diagram,
$$
d\pi_N(d\widehat{\varphi}_\xi(v))
=d\varphi_p(d\pi_M(v)).
$$
Applying the definition of
$\widehat{\varphi}(\xi)$,
we get
$$
(\widehat{\varphi}^*\lambda_N)_\xi(v)
=\xi(d\pi_M(v))=(\lambda_M)_\xi(v).
$$
The statement for $\omega_{\text{can}}$
follows by applying $-d$.
\end{proof}

\noindent
This is why cotangent lifts show up so often:
every diffeomorphism, or every action, on $M$
gives a canonical symplectic
diffeomorphism, or symplectic action,
on $T^*M$.
If $X$ is a vector field on $M$ and
$X_{T^*M}$ is the infinitesimal vector
field of the lifted flow on $T^*M$, then
$$
i_{X_{T^*M}}\omega_{\text{can}}
=d\left(\lambda(X_{T^*M})\right),
$$
so the cotangent lift is Hamiltonian with
Hamiltonian function
$$
H_X(\xi)=\xi(X_{\pi(\xi)}).
$$

\begin{exercise}
\label{exercise-cotangent-lift-char}
Let $F:T^*M \to T^*M$ be a
diffeomorphism such that
$$
F^*\lambda=\lambda.
$$
Show that $F$ is the cotangent lift of a
diffeomorphism $\varphi:M \to M$.

\noindent
Hint: first show that $F$ preserves the
vertical distribution
$\ker d\pi \subset T(T^*M)$.
Then show that $F$ sends fibers of
$T^*M$ to fibers.
\end{exercise}

\section{Moser's trick}
\label{section-moser-trick}

\noindent
Moser's trick is an application of the
following technical lemma:

\begin{lemma}
\label{lemma-lie-deriv-time-depen-vecto}
\begin{reference}
\cite[Proposition 22.15]{les}
\end{reference}
Let $M$ be a smooth family and $X_t$ a
one-parameter family of vector fields
on $M$, i.e. a {\it time-dependent vector
field}.
Let $A_t$ be a family of tensors indexed by
the same parameter $t$.
Then
$$
\frac{d}{dt}\Big|_{t=t_0}\varphi^* _tA_t
=\varphi^*_t\left(\mathcal{L}_{X_t}A_t+
\frac{d}{dt}\Big|_{t=t_0}A_t\right).
$$
\end{lemma}

\begin{theorem}[Moser's trick]
\label{theorem-moser-trick}
Let $M$ be a smooth manifold and $V_1, V_0$
two cohomologous volume forms.
Then there exists a diffeomorphism $\varphi:M
\to M$
such that $\varphi^*V_1=V_0$.
\end{theorem}

\begin{proof}
Suppose that $V_1-V_0=d\eta$ for some $\eta
\in \Lambda^{n-1}(M)$.
To construct a time-dependent vector field
consider the forms
$$
V_t:=tV_1+(1-t)V_0,\qquad t \in [0,1].
$$
Then we may associate a unique vector field
$X_t$ to $-\eta$
via the isomorphism $i_\bullet
V_t:\mathfrak{X}(M) \to \Lambda^{n-1}(M)$.
This is a time-dependent vector field,
so that by Lemma
\ref{lemma-lie-deriv-time-depen-vecto}
we see that the pullback forms
$\varphi^*_tV_t$ are constant:

\begin{align*}
\frac{d}{dt}\Big|_{t=t_0}\varphi^*_tV_t
&=\varphi^*_t
\left(\mathcal{L}_{X_t}V_t
+\frac{d}{dt}\Big|_{t=t_0}V_t\right)\\
&=\varphi^*_t
\left(di_{X_t}V_t+i_{X_t}dV_t
+\frac{d}{dt}\Big|_{t=t_0}(tV_1+(1-t)V_0)
\right)\\
&=\varphi^*_t
\left(d(-\eta)
+V_1-V_0\right)\\
&=\varphi^*_t(-(V_1-V_0)+V_1-V_0)\\
&=0.
\end{align*}

\noindent
The results follows by taking $t=1$.
\end{proof}

\section{Darboux theorem}
\label{section-darboux-theorem}

\begin{theorem}[Darboux]
\label{theorem-darboux}
Let $(M,\omega)$ be a symplectic manifold and
$p \in M$.
There exist local coordinates
$x^1,…,x^n,y^1,…,y^n$ at $p$ such that
$$
\omega=\sum_{i=1}^n dx^i \wedge dy^i.
$$
\end{theorem}

\begin{proof}
Take any chart $(\varphi,U)$ at $p$.
It's enough to show that there is a
symplectomorphism
$\psi:(U,\varphi^*\omega)\to(U,
\omega_{\text{can}})$.

Since both $\varphi^*\omega$ and
$\omega_{\text{can}}$
are closed, so is their difference,
so that by Poincaré lemma, restricting the
neighbourhood if necessary,
we may assume that
$d\eta=\varphi^*\omega-\omega_{\text{can}}$
for some 1-form $\eta$.
The result follows by applying the exact same
procedure as in
Moser's trick as follows.
Let $\omega_t:=t
\varphi^*\omega+(1-t)\omega_{\text{can}}$.
Then $\omega_t$ is symplectic for all $t$.
Assign to every $t$ the vector field
$X_t:=\omega_t^\sharp(-\eta)$.
Then apply Lemma
\ref{lemma-lie-deriv-time-depen-vecto}.
\end{proof}

\section{Weinstein's Lagrangian submanifold
theorem}
\label{section-weins-lagra-subma}

\noindent
It was Weinstein who said that everything is
a Lagrangian submanifold.
The following theorem says that Lagrangian
submanifolds are
defined intrinsically in the sense that they
can all
be thought of as embedded in its own
cotangent bundle.

\begin{theorem}[Weinstein's Lagrangian
submanifold theorem]
\label{theorem-weins-lagra-subma}
If $L \subset M$ is a Lagrangian submanifold,
there is are neighbourhoods $L \subseteq U_0
\subset M$
and $L \subset U_1 \subset T^*L$
and a symplectomorphism $\varphi:U_0 \to
U_1$.
\end{theorem}

\section{Coadjoint orbits and KKS theorem}
\label{section-coadjoint-orbits}

\noindent
The coadjoint action induces orbits in
$\mathfrak{g}^*$. The tangent space at any
point of a coadjoint orbit is given by
infinitesimal vector fields, that is, $$
T_{\xi}\mathcal{O}=\{u_{\mathfrak{g}^*}(\xi):
\forall
u \in \mathfrak{g}\}. $$

\noindent
We may introduce a symplectic form
$\omega \in \Lambda^{2}T^*_\xi\mathcal{O}$
on the tangent space
at some point of a
coadjoint orbit by
$$
\omega_\xi(X,Y):=\xi([u,v])
=-u_{\mathfrak{g}^*}(\xi)v
=v_{\mathfrak{g}^*}(\xi)(u)
$$
which does not depend on
$u,v \in \mathfrak{g}$.

\begin{theorem}[KKS]
\label{theorem-KKS}
$\xi \mapsto  \omega_\xi$ is symplectic on
$\mathcal{O}$.
\end{theorem}

\begin{proof}
Non degeneracy is easy.
Then show that $\omega$ is $G$-invariant.
Closedness follows from Jacobi identity.
\end{proof}

\begin{exercise}
\label{exercise-rough-vector-field}
Let $G$ be a Lie group.
Let $X:G \to TG$ be a section
of the projection $TG \to G$,
not necessarily smooth.
Show that if $X$ is left invariant
(i.e. $dL_g(X)=X \circ L_g$
for all $g \in G$, then $X$ is smooth.
\end{exercise}

\begin{proof}
We have to show that
$X(f)$ is a smooth function on $M$
for any $f \in C^\infty(M)$.
Indeed,
\begin{align*}
X(f)g&=X_g(f)\\
&=X_e \circ L_g(f)\\
&=dL_g(X_e)(f)\\
&=X_e(f \circ L_g),
\end{align*}

\noindent
which is smooth since $X_e$
is not varying while
$f \circ L_g$ is smooth.
\end{proof}

\begin{exercise}
\label{exercise-killing-form}
For a Lie algebra $\mathfrak{g}$,
there is always a canonical bilinear form
$k:\mathfrak{g}\times\mathfrak{g}
\to \mathbb{R}$
called the {\it Killing form}
given by
$$
k(u,v)=\text{tr}(\text{ad}_u\text{ad}_v).
$$
\begin{enumerate}
\item Note that $k$ is symmetric
and check that it is Ad-invariant.
\item A Lie algebra is called semisimple
if $k$ is nondegenerate.
Show that $\mathfrak{so}(3)$ is semisimple.
\end{enumerate}
\end{exercise}

\begin{proof}
\begin{enumerate}
\item
\item Let
$$
e_1=\begin{pmatrix}
0&1&0\\
-1&0&0\\
0&0&0\end{pmatrix},
e_2=\begin{pmatrix}
0&0&1\\
0&0&0\\
-1&0&0\end{pmatrix},
e_3=\begin{pmatrix}
0&0&0\\
0&0&1\\
0&-1&0\end{pmatrix}
$$
be a basis of $\mathfrak{so}_3$.
Then
$$
[e_1,e_2]=-e_3,\quad
[e_1,e_3]=e_2,\quad
[e_2,e_3]=-e_1.
$$
To compute the trace we use
Euclidean dot product identifying
$\mathfrak{so}_3=\mathbb{R}^3$.
We obtain
\begin{align*}
\kappa(e_1,e_2)&=
\text{tr}\text{ad}_{e_1}\text{ad}_{e_2}
&=\sum_i\left<[e_1,[e_2,e_i]],e_i\right>\\
&=\left<[e_1,[e_2,e_1]],e_1\right>
+\left<[e_1,[e_2,e_2]],e_2\right>
+\left<[e_1,[e_2,e_3]],e_3\right>\\
&=0,
\end{align*}

\noindent
and similarly
$$
\kappa(e_1,e_3)=1,\quad
\kappa(e_2,e_1)=1,\quad
\kappa(e_3,e_3)=-2.
$$
Ok, the computations are left to the reader,
but $\kappa$ is should be a diagonal matrix
and therefore nondegenerate.
\end{enumerate}
\end{proof}
\section{Hamiltonian actions and moment map}
\label{section-hamiltonian-action-moment-map}

\noindent
To warm-up we define real symplectic and
Hamiltonian actions.

\begin{definition}
\label{definition-one-parameter-actions}
\begin{itemize}
\item An $\mathbb{R}$-action is {\it
symplectic}
if $\psi_t^*\omega=\omega
\iff L_X \omega=0$,
that is, $X$ is a symplectic field.
\item An $\mathbb{R}$-action is {\it
hamiltonian}
if $X$ is Hamiltonian,
that is, if there exists
a function $H \in C^\infty(M)$ such that
$X=X_H$.
\end{itemize}
\end{definition}

\noindent
To define moment map suppose
there exists a map $\hat{\mu}$
such that
$$
\xymatrixcolsep{3em}
\xymatrixrowsep{.7em}
\xymatrix{
&C^\infty(M)\ar[dd]&H\ar@{|->}[dd]\\
\mathfrak{g}\ar[dr]\ar[ur]^{\hat{\mu}}\\
u \ar@{|->}[dr] & \mathfrak{X}(M) & X_H\\
& u_M=X_u \ar@{=}[ur]
}
$$
Note that if  $\hat{\mu}$ exists, it's not
uniquely determined, and that we may suppose
that it is linear.

It's equivalent to define $\hat{\mu}$
as a map
$\hat{\mu}:\mathfrak{g}\to C^\infty(M)$
or as a map
\begin{align*}
\mu: M &\longrightarrow \mathfrak{g}^* \\
\hat{\mu}(u)(x) &= \left<\mu(x),u\right>.
\end{align*}

\begin{definition}
\label{definition-moment-map}
\begin{itemize}
\item A $G$-action is symplectic if
$\psi_g^*\omega=\omega$ for all $g \in G$,
that is, if
$$
\xymatrix{
\psi:G\ar[r]\ar[dr]&\text{Dif}(M)\\
&\text{Symp}(M,\omega)\ar@{^{(}->}[u].
}
$$

\item A symplectic action of $G$ on
$(M,\omega)$ via $\psi$
is {\it weakly Hamiltonian}
if there exists a map
$\mu:M \to \mathfrak{g}^*$
(or equivalently,
$\hat{\mu}:\mathfrak{g} \to C^\infty(M)$
linear) such that
$$
i_{u_M}\omega=d\left<\mu(p),u\right>,\qquad
(\text{resp. } u_M=X_{\hat{\mu}(u)}).
$$

\item A weakly Hamiltonian action is
{\it Hamiltonian} if it is also
$\mu$-equivariant for the action of $G$ on
$M$ and the coadjoint action on
$\mathfrak{g}^*$, that is,
$$
\mu \circ \psi_g=(\text{Ad}^*)_g(\mu).
$$
In this case we call $\mu$ a {\it moment
map}.
\end{itemize}
\end{definition}

\begin{exercise}
\label{exercise-linear-moment-map-rotation}
Let $V=\mathbb{R}^2$ with symplectic form
$\omega=dx\wedge dy$.
Consider the action of $S^1$ on $V$
given by rotations,
$$
e^{it}(x,y)=
(x\cos t-y\sin t,x\sin t+y\cos t).
$$
Identifying $\mathfrak{s}^1=\mathbb{R}$,
let $u=1$.
\begin{enumerate}
\item Compute the infinitesimal vector field
$$
u_V(x,y)=
\frac{d}{dt}\Big|_{t=0} e^{it}(x,y).
$$
\item Compute the $1$-form $i_{u_V}\omega$.
\item Find a function $H:\mathbb{R}^2
\to \mathbb{R}$ such that
$$
dH=i_{u_V}\omega.
$$
\item Check that
$$
H(v)=\frac{1}{2}\omega(uv,v).
$$
Here we identify $T_vV$ with $V$, and $uv$
denotes the infinitesimal action of $u$ on
$v$.
\end{enumerate}
\end{exercise}

\begin{proof}
We have
\begin{align*}
u_V(x,y)
&=
\frac{d}{dt}\Big|_{t=0}
(x\cos t-y\sin t,
x\sin t+y\cos t)\\
&=(-y,x).
\end{align*}
Let $w=(a,b)$. Since
$$
\omega((p,q),(a,b))=pb-qa,
$$
we get
\begin{align*}
(i_{u_V}\omega)(w)
&=\omega((-y,x),(a,b))\\
&=-yb-xa.
\end{align*}
Thus
$$
i_{u_V}\omega=-x\,dx-y\,dy.
$$
Therefore
$$
H(x,y)=-\frac{1}{2}(x^2+y^2)
$$
satisfies
$$
dH=i_{u_V}\omega.
$$
Finally, we identify $T_vV$ with $V$. Since
$$
uv=(-y,x),\qquad v=(x,y),
$$
we get
\begin{align*}
\frac{1}{2}\omega(uv,v)
&=
\frac{1}{2}\omega((-y,x),(x,y))\\
&=
-\frac{1}{2}(x^2+y^2).
\end{align*}
Hence
$$
H(v)=\frac{1}{2}\omega(uv,v).
$$
\end{proof}

\noindent
The following exercise shows
how the moment map, i.e.
$\hat{\mu}=H$ looks for any
linear symplectic action.

\begin{exercise}
\label{exercise-linear-moment-map}
Let $(V,\omega)$ be a symplectic vector space
and let $G$ act linearly on $V$ preserving
$\omega$.
Write
$$
u_M(v)=X_u(v)=uv=\frac{d}{dt}\Big|_{t=0}
\exp(tu)v
$$
for the infinitesimal action of $u$ on $v$.
We identify $T_vV$ with $V$ and,
for $u \in \mathfrak{g}$, define
$$
H_u(v)=\frac{1}{2}\omega(uv,v).
$$
Show that
$$
dH_u=i_{u_V}\omega,
$$
where $u_V(v)=uv$ is the infinitesimal
vector field associated to $u$.
\end{exercise}
\begin{proof}
Let $w \in V$. Since $V$ is a vector
space, we identify $T_vV$ with $V$.
Thus $w$ is a tangent vector at $v$,
and $dH_u(v)(w)$ is the directional
derivative of $H_u$ at $v$ in the
direction $w$. Therefore
\begin{align*}
dH_u(v)(w)
&=
\frac{d}{dt}\Big|_{t=0}
H_u(v+tw)\\
&=
\frac{1}{2}
\frac{d}{dt}\Big|_{t=0}
\omega(u(v+tw),v+tw)\\
&=
\frac{1}{2}\omega(uw,v)
+
\frac{1}{2}\omega(uv,w).
\end{align*}
Here we used that $\omega$ is bilinear:
differentiating a bilinear expression
means differentiating in the first slot
and in the second slot.

Since the action preserves $\omega$,
we have
$$
\omega(e^{tu}a,e^{tu}b)=\omega(a,b)
$$
for all $a,b \in V$.
Differentiating at $t=0$ gives
$$
\omega(ua,b)+\omega(a,ub)=0.
$$
Taking $a=w$ and $b=v$, we obtain
$$
\omega(uw,v)+\omega(w,uv)=0.
$$
Since $\omega$ is skew-symmetric, this
gives
$$
\omega(uw,v)=\omega(uv,w).
$$
Therefore
$$
dH_u(v)(w)=\omega(uv,w).
$$
But $u_V(v)=uv$, so
$$
dH_u(v)(w)
=
(i_{u_V}\omega)_v(w).
$$
Thus
$$
dH_u=i_{u_V}\omega.
$$
\end{proof}

\begin{exercise}
\label{exercise-tangent-moment-level-set}
Let $(M,\omega)$ be a symplectic manifold
and let $G$ act on $M$ with moment map
$\mu:M\to\mathfrak{g}^*$.
Assume that $0$ is a regular value of
$\mu$.
Show that, for every $p \in \mu^{-1}(0)$,
we have
$$
T_p\mu^{-1}(0)
=
(T_p(Gp))^\omega.
$$
Here $(T_p(Gp))^\omega$ denotes the
symplectic orthogonal of the tangent
space to the orbit.
\end{exercise}

\begin{proof}
For the inclusion $\subseteq$
we need to show that $\omega(v,w)=0$
for any $v \in T_p\mu^{-1}(0)$
and $w \in T_p(G\cdot p)$.
A $w$ looks like: take a 1-parameter
subgroup of $G$, which of course
we want to see as something of the kind
$\text{exp}(tu)\cdot p$ and differentiate
at $t=0$. And of course this is the
infinitesimal generator $u_M = X_u$
for $u \in \mathfrak{g}$.
On the other hand, a $v$ looks like:
take a curve $\gamma \subset \mu^{-1}(0)$
and differentiate at $t=0$.

By the moment map property we see that
$$
\omega(u_M,v)=d\left<\mu,u\right>_p(v).
$$
How to compute that differential?
I remember that the differential of a
function can be computed by
using a curve that realises the vector,
i.e. $df(v) = \frac{d}{dt}\Big|_{t=0}
f\circ \gamma(t)$.
That is,
$$
d\left<\mu(p),u\right>v
=\frac{d}{dt}\Big|_{t=0}
\left<\mu(\gamma(t)),u\right>=0.
$$

For the $\supseteq$ inclusion
suppose that $v$ is such that
$\omega(u_M,v)=0$ for all
$u \in \mathfrak{g}$.
To show that $v \in T_p\mu^{-1}(0)$
we need to show that
$v \in \Ker d\mu$ by the
level set theorem.
This is what the moment map equation says:
$0=i_{u_M}\omega(v)
=d\left<\mu,u\right>_p(v)$,
and the latter differential is in fact
$d(\text{ev}_u \circ \mu)$.
Since $\operatorname{ev}_u$ is linear, its
differential is again $\operatorname{ev}_u$.
Hence
$$
d\left<\mu,u\right>_p(v)
=
\left<d\mu_p(v),u\right>.
$$
Since this vanishes for every
$u \in \mathfrak{g}$, the covector
$d\mu_p(v)\in\mathfrak{g}^*$ vanishes
identically. Thus $d\mu_p(v)=0$.
\end{proof}

\begin{exercise}
\label{exercise-reduced-symplectic-form}
Let $(M,\omega)$ be a symplectic manifold
with a Hamiltonian action of $G$ and
moment map $\mu:M\to\mathfrak{g}^*$.
Assume that $0$ is a regular value of
$\mu$ and that $G$ acts freely and
properly on $\mu^{-1}(0)$.

Let
$$
Z=\mu^{-1}(0)
$$
and let
$$
\pi:Z\longrightarrow Z/G
$$
be the quotient map.
Show that there exists a unique $2$-form
$\omega_{\mathrm{red}}$ on $Z/G$ such that
$$
\pi^*\omega_{\mathrm{red}}
=
\iota^*\omega,
$$
where $\iota:Z\hookrightarrow M$ is the
inclusion.
\end{exercise}

\begin{proof}
By Smooth Manifolds Exercise
\ref{exercise-basic-form},
we know that all we have to do is check that
$\Ker \pi_* \subset \Ker\omega$,
that is, that every vertical vector field
$Z$ gives the zero 1-form when contracted
with $\omega$, i.e. $i_Z\omega=0$.

We know that every vertical vector field
is locally given as $u_Z$ for some
$u \in \mathfrak{g}$.
By the moment map equation,
we shall obtain that $i_{u_Z}\omega=0$
if we show that $\left<\mu(\cdot),u\right>$
is a constant function on $Z$,
which is clear since $Z= \mu^{-1}(0)$.
\end{proof}

\noindent
Note that the original intention of the
exercise above was to use its
preceding exercise (but of course
my solution is also correct).
Here's the alternative proof:

By Exercise \ref{exercise-tangent-moment-level-set},
we have
$$
T_pZ=(T_p(G\cdot p))^\omega.
$$
Let $X \in \Ker \pi_{*,p}$ and
$Y \in T_pZ$. Since $X$ is vertical,
we have
$$
X \in T_p(G\cdot p).
$$
Since $Y \in T_pZ$, the previous identity
gives
$$
Y \in (T_p(G\cdot p))^\omega.
$$
Hence
$$
(\iota^*\omega)_p(X,Y)=\omega_p(X,Y)=0.
$$
Thus
$$
\Ker\pi_* \subset \Ker(\iota^*\omega).
$$
Since $\iota^*\omega$ is closed, Smooth
Manifolds Exercise \ref{exercise-basic-form}
implies that $\iota^*\omega$ is basic.
Therefore it descends to a unique $2$-form
$\omega_{\mathrm{red}}$ on $Z/G$.



\begin{exercise}
\label{exercise-reduced-form-is-symplectic}
In the setting of Exercise
\ref{exercise-reduced-symplectic-form},
show that the reduced form
$\omega_{\mathrm{red}}$ is symplectic.
That is, show that
$$
d\omega_{\mathrm{red}}=0
$$
and that $\omega_{\mathrm{red}}$ is
nondegenerate.
\end{exercise}

\begin{exercise}
\label{exercise-reduction-cpn}
This result is the geometric engine behind symplectic reduction. Let's see it in action for the most important example: the construction of $\mathbb{CP}^{n-1}$ from $\mathbb{C}^n$.
Consider $M = \mathbb{C}^n \cong \mathbb{R}^{2n}$ with the standard symplectic form $\omega = \sum dx_j \wedge dy_j$. Let $S^1$ act by diagonal rotation: $e^{i\theta} \cdot z = e^{i\theta}z$.
\begin{enumerate}
\item Show that the moment map for this action is $\mu(z) = \frac{1}{2}\sum |z_j|^2 + c$. We choose $c = -1/2$ so that $\mu^{-1}(0) = S^{2n-1}$ is the unit sphere.
\item Show that the orbits of this action on $S^{2n-1}$ are the fibers of the Hopf fibration.
\item Using the previous exercise, check that $T_p S^{2n-1} = (T_p(\text{Orbit}))^\omega$. This means the restricted form $\omega|_{S^{2n-1}}$ is degenerate only in the direction of the orbits!
\item This implies that $\omega$ descends to a non-degenerate form on $S^{2n-1}/S^1 = \mathbb{CP}^{n-1}$. This is (a multiple of) the Fubini-Study form.
\end{enumerate}
\end{exercise}

\section{Noether's principle}
\label{section-Noether-principle}

\begin{theorem}[Noether's principle]
\label{theorem-Noether-principle}
Let $(M,\omega,\mu)$ be a Hamiltonian
$G$-space with Hamiltonian
$H\in C^\infty(M)$. $H$ is $G$-invariant if
and only if $\mu$ is preserved by
the Hamiltonian flow.
\end{theorem}

\begin{proof}
one line
\end{proof}

\section{Symplectic quotient}
\label{section-symplectic-quotient}

\noindent
In symplectic geometry we cannot always take
quotients.
We need to use level-sets of the moment map.

\begin{theorem}[Marsden-Weinstein, Megre]
\label{theorem-marsden-weinstein-megre}
Let $G$ act on $(M,\omega)$ in a Hamiltonian
way
with moment map $\mu:M \to \mathfrak{g}^*$.
Suppose $0 \in \mathfrak{g}^*$ is a regular
value.
By the equivariance of $\mu$, we have an
action of $G$ on $\mu^{-1}(0)$
given as follows. For $x \in \mu^{-1}(0)$,
$\mu(\psi_g(x))=\text{Ad}_g^*(\mu(x))=0$.
Suppose that this action is regular (this
will be defined later),
that $\mu^{-1}(0)/G$ is smooth, and that the
projection is a submersion.
Then there exists a unique symplectic form
$\omega_{\text{red}}$ in $M_0$
such that
$i^*\omega=\pi_0^*\omega_{\text{red}}$.
\end{theorem}

\noindent
That is, the pullback of the symplectic form
to the
0 level-set of $\mu$ passes to the quotient:
$$
\xymatrix{
G\text{ acts on
}\mu^{-1}(0)\ar[d]^{\pi_0}\ar@{^{(}->}[r]&
M\\
(M_0:=\mu^{-1}(0)/G,\omega_{\text{red}})
}
$$

\noindent
A slightly more general setting is the
following.
Hamiltonian action of $G$ on $(M,\omega)$,
and moment map
$\mu:M \to \mathfrak{g}^*$.
$\xi \in \mathfrak{g}^*$ regular, $G_\xi
\subseteq G$
stabilizer of $\xi$.
Then we have
$$
\xymatrix{
G_\xi\text{ acts on
}\mu^{-1}(\xi)\ar[d]^{\pi_\xi}\ar@{^{(}->}[r]
^-{i_\xi}&
M\\
(M_\xi=\mu^{-1}(\xi)/G_\xi,\omega_{\text{red}
}).
}
$$

\begin{exercise}
\label{exercise-so3-action-cotangent}
Consider the group $\text{SO}(3)$ acting
on $T^*\mathbb{R}^{3}$ by the cotangent
lift of the usual action of
$\text{SO}(3)$ on $\mathbb{R}^{3}$.
\begin{enumerate}
\item For $u\in\mathfrak{so}(3)$, compute
the corresponding infinitesimal generator
$u_{T^*\mathbb{R}^3}\in\mathfrak{X}
(T^*\mathbb{R}^{3})$.
\item Identify $\mathfrak{so}(3)$ with
$\mathbb{R}^3$ as in Lista 5. Show that,
with this identification, we have
$u_{T^*\mathbb{R}^3}(q,p)
=(u\times q,u\times p)$.
\item Identifying $\mathfrak{so}(3)^*
\cong(\mathbb{R}^{3})^*
\cong\mathbb{R}^{3}$ using the usual
inner product, show that the moment map
for the action of $\text{SO}(3)$ on
$T^*\mathbb{R}^{3}$ is $\mu(q,p)=q\times
p$. Conclude (by Noether's theorem) that
if $V\in \mathcal{C}^\infty(\mathbb{R}^3)$
is $\text{SO}(3)$-invariant, then the
flow of the Hamiltonian $H(q,p)=\frac{p^2}
{2m}+V(q)$ preserves ``angular momentum''
$q \times p$.
\end{enumerate}
\end{exercise}

\begin{proof}
\begin{enumerate}
\item To survive without getting too
involved in the definitions, I observe
the following. $\text{SO}(3)$ acts
by rotations on $\mathbb{R}^3$.
This means that the tangent vector
to any 1-parameter orbit is a tangent
vector to a sphere containing the
1-parameter orbit.
In other words, we have an expression
for the 1-parameter orbit
$p \cos t + u_{\mathbb{R}^3} \sin t$
where $u_{\mathbb{R}^3}$
is the $\mathbb{R}^3$ component
of $u_{T^*\mathbb{R}^3}$ at
$p \in \mathbb{R}^3$.
This just says that $u_{\mathbb{R}^3}$
is orthogonal to $p$.
Since the induced map on the cotangent
space is essentially given by the
pushforward of the action, which is linear,
then it's exactly the same ---
thus the $T^*\mathbb{R}^3$
component of $u_{T^*\mathbb{R}^3}$
must also be orthogonal to
the covector where it is based.

Now let's get too involved in the
definitions.
By definition, for
$(p,q) \in T^*\mathbb{R}^3$
and $u \in \mathfrak{so}_3$,
$$
u_{T^*\mathbb{R}^3}
=\frac{d}{dt}\Big|_{t=0}
e^{tu}(p,q).
$$
To find the integral curve $\varphi(t)
=e^{tu}$ recall how it is constructed:
the vector $u \in T_e G$
gives the left-invariant vector field
$X(g)=d_eL_g(u)$, which
has an integral curve at the identity
$\varphi_u(t)$ defined by the properties
$\varphi_u'(0)=u$ and $\varphi_u(0)=e$.
We denote $\varphi_u(t)=e^{tu}$.
Thus differentiating the integral
curve by definition gives
the left-invariant vector field $X$
along $\varphi_u$,
which is the differential of a linear
map and thus it acts by itself,
that is $X(\varphi(t))
=d_eL_{\varphi_u(t)}(u)
=\varphi_u(t)u$.

Differentiating and evaluating at $t=0$
we see that $u_{T^*\mathbb{R}^3}
=(up,uq)$
since, in fact, the action on the
cotangent part is also given by
simple matrix multiplication.

More explicitly, let $A\in \text{SO}(3)$.
The cotangent lift is encoded by
$$
\xymatrix{
T^*\mathbb{R}^3 \ar[r]^{\widehat A}
\ar[d]_{\pi}
& T^*\mathbb{R}^3 \ar[d]^{\pi}\\
\mathbb{R}^3 \ar[r]_{A}
& \mathbb{R}^3
}
$$
and sends $(q,\xi)$ to $(Aq,A\cdot \xi)$,
where
$$
(A\cdot \xi)(v)=\xi(A^{-1}v).
$$
Use the Euclidean inner product to
identify covectors with vectors. Thus
write
$$
\xi(v)=\langle p,v\rangle.
$$
Here $p$ is the vector representing
$\xi$ by the Riesz representation
theorem.
By the definition of adjoint from
Definition \ref{definition-adjoint},
we have
$$
\langle A^{-1}v,p\rangle
=\langle v,(A^{-1})^\dag p\rangle.
$$
Therefore
\begin{align*}
(A\cdot \xi)(v)
&=\xi(A^{-1}v)\\
&=\langle p,A^{-1}v\rangle\\
&=\langle A^{-1}v,p\rangle\\
&=\langle v,(A^{-1})^\dag p\rangle\\
&=\langle (A^{-1})^\dag p,v\rangle.
\end{align*}
Thus the covector $A\cdot \xi$ is
represented by $(A^{-1})^\dag p$.
Since $A\in\text{SO}(3)$, we have
$A^\dag=A^{-1}$, and hence
$(A^{-1})^\dag=A$. Therefore, after this
identification,
$$
A\cdot(q,p)=(Aq,Ap).
$$
It follows that, for
$u\in\mathfrak{so}(3)$,
$$
u_{T^*\mathbb{R}^3}(q,p)
=
\frac{d}{dt}\Big|_{t=0}
(e^{tu}q,e^{tu}p)
=(uq,up).
$$
\item We now identify $u$ with a vector
and compute $Uq$. Write
$$
U=
\begin{pmatrix}
0&-c&b\\
c&0&-a\\
-b&a&0
\end{pmatrix}
\longleftrightarrow (a,b,c).
$$
If $q=(q_1,q_2,q_3)$, then
\begin{align*}
Uq
&=
\begin{pmatrix}
0&-c&b\\
c&0&-a\\
-b&a&0
\end{pmatrix}
\begin{pmatrix}
q_1\\ q_2\\ q_3
\end{pmatrix}\\
&=
\begin{pmatrix}
-cq_2+bq_3\\
cq_1-aq_3\\
-bq_1+aq_2
\end{pmatrix}.
\end{align*}
These are the coordinates of
$$
(a,b,c)\times(q_1,q_2,q_3).
$$
Thus $Uq=u\times q$. The same computation
for $p$ gives $Up=u\times p$. Therefore
$$
u_{T^*\mathbb{R}^3}(q,p)
=(u\times q,u\times p).
$$

\item There is of course the tour-de-force
path in which we do everything in
coordinates and just confirm that
the moment map and equivariance equations
hold. But let's not do that.

I see two other ways: (1) try to use somehow
that the symplectic form in $T^*\mathbb{R}^3$
is the differential of the tautological
form, and (2) try to use that everything
is linear.

Let's do (2). Recall that $\mu$
takes values in dual Lie algebra.
Thus the vector $q\times p$ we are given
is to be understood as a traceless
anitymmetric matrix. The map
$T^*\mathbb{R}^3 \ni (q,p) 
\overset{\hat{\mu}}{\mapsto } 
(q\times p)u$
is most likely linear.
So the exterior differential is…
probably itself!
\end{enumerate}
\end{proof}

\section{Toric varieties}
\label{section-toric-varieties}

\noindent
For any compact connected symplectic manifold
$(M^{2n},\omega)$ with an action
of a torus $\mathbb{T}^n$ and moment map
$\mu:M \to \mathbb{R}^n
=(\mathfrak{t}^n)^*$, the image of $\mu$
is a polytope. By [Guilleimin-Sternberg] and
[Atiyah], this polytope is
convex. Further, if the dimension of the
torus is exactly half
the dimension of the manifold, by [Delzant],
the moment polytope
determines the manifold, i.e. there is a
one-to-one
correspondence between
Delzant polytopes and symplectic toric
varieties.
Moduli space of polygons is an example of
this,
see \cite{moduli-polygons}.

In fact, this situation of having a torus
action of half the dimension
is what we may call a ``totally integrable
system''.

\section{(Completely) integrable systems}
\label{section-completely-integrable-systems}

\noindent
Let $(M,\omega)$ be a symplectic manifold.
Recall that

\begin{itemize}
\item
$\{f,g\}=\omega(X_g,X_f)=L_{X_f}g=-L_{X_g}f$.
\item $X_{\{f,g\}}=[X_f,X_g]$
\end{itemize}

\noindent
If $\{f,g\}=0$ then the Hamiltonian flows of
$f$ and $g$ commute.

\begin{definition}
\label{definition-compl-integ-syste}
Let $(M^{2n},\omega)$ and $H \in
C^\infty(M)$.
$(M,\omega,H)$ is a {\it completely
integrable system}
(in Liouville's sense)
if there exist functions $f_1,\ldots,f_n \in
C^\infty(M)$ such that
\begin{enumerate}
\item $H=f_1,\ldots,f_n$ are linearly
independent
($df_1,\ldots,df_n$ are linearly independent
in every point).

\item $\{ f_i, f_j\}=0$ for all $i,j$.
\end{enumerate}
\end{definition}

\noindent
Note that $F=(f_1,\ldots,f_n):M \to
\mathbb{R}^n$ is a submersion.
For all $c \in \mathbb{R}^n$, $F^{-1}(c)$
are Lagrangian submanifolds.
$F^{-1}(c)$ is invariant by the flow of
$X_{f_i}$ for $i=1,\ldots,n$.

Explanation: $F^{-1}(c)$ is a level-set
submanifold of $M$,
whose tangent space is generated by the
$df_i$.
These are all in the kernel of $dF$ by the
integrability condition,
namely $dF=(df_1,\ldots,df_n)$,
$df_i(X_{f_j})=\{f_i,f_j\}=0$.

The idea is that finding these ``conserved
quantities''
we manage to describe the Hamiltonian
dynamics completely.
Showing that a system is integrable is a
victory.
The moduli space of polygons is an example,
and I think that so is
$(\mathbb{C}^2)^{[n]}$.

There is another method of using
``bihamiltonian actions''
(two Hamiltonian actions yielding the same
dynamics)
/Lax pairs (defining the system using
eigenvalues of matrices).
Examples of this are Toda, Caolengo-Moser.

\medskip\noindent
The following theorem says that the connected
component of a symplectic manifold with an
integrable system


\begin{theorem}[Arnold-Liouville]
\label{theorem-arnold-liouville}
Let $(M^{2n},\omega)$, $H=f_1,\ldots,f_n$ be
completely intgrable system
and $F:M \to \mathbb{R}^n$,
$F=(f_1,\ldots,f_n)$.
Let $M_c$ be the connected component of
$F^{-1}(c)$.
Then
\begin{enumerate}
\item
\label{item-toric1}
If $M_c$ is compact, then $M_c \simeq
\mathbb{T}^n=\mathbb{R}^n/2\pi\mathbb{Z}^n=
\{(\theta_1,\ldots,\theta_n)
\text{
mod }2\pi\}$.

\item
\label{item-toric2}
In angular coordinates, the Hamiltonian flow
is linear, that is,
$\frac{d}{dt}^{-1}(t)=\alpha \in
\mathbb{R}^n$
implies that $\theta(t)=\theta_0+\alpha t$.
\end{enumerate}
\end{theorem}

The proof of item \ref{item-toric1} is given
by
the following proposition
putting $N=M$ and
$X_1=X_{f_1},\ldots,X_n=X_{f_n}$.

\begin{proposition}
\label{proposition-torus}
Let $N$ be complete, connected, of dimension
$n$.
Let $X_1,\ldots,X_n$ be linearly independent
vector fields on $N$
such that $[X_i,X_j]=0$ for all $j$.
Then $N \simeq \mathbb{T}^n$.
\end{proposition}

\begin{proof}
First we present the steps:
\begin{enumerate}
\item Flows of $X_1,\ldots,X_n$ define an
action of $\mathbb{R}^n$ on $N$.

\item This action is transitive. Then $N
\simeq \mathbb{R}^n/\Gamma$
where $\Gamma$ is the stabilizer.

\item Since $\dim N=n$, then  $\Gamma$ is a
lattice of $\mathbb{R}^n$.

\item Since $N$ is compact, then $\Gamma$ is
of maximum rank, that is,
$\Gamma=2\pi \mathbb{Z}^n \subset
\mathbb{R}^n$ .
\end{enumerate}


Let $\varphi^i_t$ be the flow of $X_i$ (which
are complete since $N$ is
compact).
Since $[X_i,X_j]=0$, then $\varphi_t^i
\varphi_s^j=\varphi_s^j \varphi^i_t$.
Thus we have an action of $\mathbb{R}^n$ on
$N$
given by moving along the flows of the vector
fields, that is
$$
\underbrace{(t_1,\ldots,t_n)}_{\mathbf{t}}
\cdot
x
=\underbrace{\varphi_{t_1}^1 \circ \ldots
\circ \varphi_{t_n}^n(x)}_{
\varphi_{\mathbf{t}}(x)}
$$
To prove transitivity
we can show that the orbit maps of this
action are surjective.
Indeed, the image of the orbit map is open
since it is a local
diffeomorphism (it's derivative is given by
the vector fields,
which are linearly independent!),
and it is closed by a taking a point in the
boundary and
using the flow at this point (and
commutativity of the flows).

To check that the stabilizer $\Gamma$ of this
action is trivial
first note that all stabilizers are equal
(not only conjugate) since $\mathbb{R}^n$ is
abelian.
Further, since the orbit map is a local
diffeomorphism,
$\Gamma$ is discrete. Thus, it must be a
lattice of rank $k \leq n$.

To show that in fact $k=n$ consider a linear
map  $T$
that maps the generators of $\Gamma$ to the
first $k$ canonical
vectors of $\mathbb{R}^n$. By the following
diagram, the arrow on the bottom
is well defined as in fact a bijective local
diffeomorphism,
that is, a diffeomorphism.
$$
\xymatrix{
\mathbb{R}^k \times\mathbb{R}^{n-k} \simeq
\mathbb{R}^n\ar[r]^{T}\ar[d]_{p}
\ar[dr]^{\psi_x \circ T}&
\mathbb{R}^n\ar[d]^{\psi_x=
\substack{\text{orbit}
\\ \text{map}}}\\
\mathbb{T}^k \times
\mathbb{R}^{n-k}\ar[r]_{\overline{T}}&
N
}
$$
Thus the factor $\mathbb{R}^{n-k}$ must be a
point since $N$ is
compact by hypothesis.
\end{proof}

\noindent
Item \ref{item-toric2} from Theorem
\ref{theorem-arnold-liouville}
is obtained by considering the flows
on the latter diagram and is left as an
exercise.

\noindent
In the case that $N$ is not compact we just
conclude that
it $M_c$ is a product of a torus with an
Euclidean space.

Perhaps the simplest example of an integrable
system
is given by taking an open ball about the
origin $B \subseteq \mathbb{R}^n$
and $B \times
\mathbb{T}^n=\{(I_1,\ldots,I_n,\phi_1,\ldots,
\phi_n)\}$
with the symplectic form $\omega=\sum_{i=1}^n
dI_1 \wedge d\phi_i$
with $\{I_1,I_j\}=0$.

The following theorem says that in fact the
{\bf symplectic form} of every
level set of the kind $M_c$, which we just
proved
that is a torus, may is described locally as
such a product.

\begin{theorem}
\label{theorem-local}
There is a neighbourhood $U$ of $M_c$ and
a symplectomorphism
$$
\xymatrix{
\sum d \mathcal{I} \wedge d\phi_i&\ar@{~>}[l]
B \times
\mathbb{T}^n\ar[r]^{r}_{\simeq}\ar[d]_{p}&U
\subseteq M\ar[d]^{F}
\ar@{~>}[r]
& \omega\\
&B\ar[r]_{\simeq}&F(U)
}
$$
\end{theorem}

\noindent
The coordinates
$(I_1,\ldots,I_n,\phi_1,\ldots,\phi_n)$
are called {\it action-angle} coordinates.

The case when $h(I)$ does not depend on the
angle coordinates $\phi_i$
we can solve the Hamiltonian equations
easily:
$$
\begin{cases}
\dot I_i=0 \\
\dot \phi_i=-\frac{\partial h}{\partial I}(I)
\end{cases}
$$
so that $\phi_i(t)=\phi_0+\alpha t$ for some
constants $\alpha$.
This is why integrable systems are actually
easy to solve!

\section{Vertex algebras}
\label{section-vertex-algebras}

Mônadas, álgebras sobre mônadas. Exemplo:
ultrafiltros (Manes, 1976). Campo
quântico, como serie de Fourier. Exemplo:
Heisenberg algebra.

Problema: o producto de campos quânticos
não
é bem definido.

\bibliography{my}
\bibliographystyle{amsalpha}

\end{document}
